% !TEX root = ../main.tex
\lecture{5}{2026-10-12}{Reele Funktionen}


\section{Reele Funktionen}
\begin{definition}
    [reele Funktion] sind Funktionen, deren Definitions- und Zielmengen Teilmengen der reelen Zahlen sind.
    \begin{equation}
        D,Z \subseteq \mathbb{R} 
    \end{equation}
    Wenn nicht explizit angegeben: natürlicher Definitionsbereich (Teilmenge der reelen Zahlen, mit Bedingung durch die Funktionslogik gegeben)
    \begin{equation}
        D_f = \left\{x \in \mathbb{R} \mid f(x)\right\}
    \end{equation}
\end{definition}

\begin{definition}
    [Wichtige Reele Funktionen]
    $n,m \in N_0$ und $a_0, \dots, a_n, b_0,\dots,b \in \mathbb{R}$
\begin{itemize}
    \item\textbf{ Affin-lineare Funktion:} $f(x) = a_0 + a_1x $
    \item \textbf{Polynom n-ten grades}
    \begin{itemize}
        \item $a_0 + a_1x^1 + \dots, a_nx^n$
    \end{itemize}
    \item \textbf{Rationale Funktionen}
    \begin{itemize}
        \item Quotient zweier Polynome
    \end{itemize}
    \item \textbf{Potenzfunktion}
    \begin{itemize}
        \item $f(x)  x^b $
        \item Definitionsbereich hängt von b ab, wenn b rationale Zahl dann Wurzel mit$ D = R \ {0}$
    \end{itemize}
    \item Exponential Funktion
    \begin{itemize}
        \item $f(x) =b^x$
    \end{itemize}
    \item natürliche Exponentialfunktion
    \begin{itemize}
        \item $f(x) = exp(x) = \textit{e}^x$
    \end{itemize}
    \item logarithmische FUnktion\begin{itemize}
        \item $f(x_) = log_b(x)$
    \end{itemize}
    \item Trigonometische Funktionen
    \begin{itemize}
        \item sin(x)
        \item cos(x)
    \end{itemize}
    \item Betragsfunktion
    \begin{itemize}
        \item $f(x) = \mid x \mid  $
    \end{itemize}
    \item Auf- und Abrundungsfunktion
    \begin{itemize}
        \item Auf: $f(x) = \lceil x \rceil = min\left\{y \in Z | y \ge x\right\}$
        \item Ab:  $f(x) = \lfloor x \rfloor = max\left\{y \in Z | y \le x\right\}$
    \end{itemize}
\end{itemize}        
\end{definition}

